\paragraph{Theorem F.2.37 (Proved: a complete canonical-representative algorithm).}
Every marked linear tuple in F.2, for either $k=2$ or $k=3$,
has a computable canonical integral representative and a
returned simultaneous unimodular conjugator. Two inputs have
the same representative if and only if they are simultaneously
$\mathrm{GL}_6(\mathbb Z)$ conjugate. The algorithm uses a proved
finite bound depending on the input. No fixed global entry bound
can represent every class.

This is an effective classification, not a claim that its
infinite output has been printed in a finite table. The
ordered involution sector has the separate complete finite
normal form F.2.38, and all of its $56$ classes have been
computed. The general bounded fallback is potentially very
expensive; its finiteness is a theorem, not a runtime estimate.

\begin{proof}
We first remove the existential condition about affine lifts
from the admissibility test for a linear tuple. A finite-order
matrix $M$ admits a free finite affine lift on $T^6$ if and
only if $\ker(M-1)\ne0$. If this kernel is zero, $1-M$ is a
surjective endomorphism of the torus, so every affine map with
linear part $M$ has a fixed point. Conversely let the order of
$M$ be $d>1$ and choose a primitive $u\in\ker(M-1)\cap\mathbb Z^6$.
The map $x\mapsto Mx+u/d^2$ has actual order $d^2$: its
$d$th power is translation by $u/d$, of exact order $d$.
For each prime $p\mid d$, the generator of its order-$p$
subgroup is translation by $u/p$, which has no fixed point.
Every nontrivial stabilizer contains a prime-order subgroup,
so the whole cyclic action is free. If $M=1$, translation
by a primitive half-vector is a free lift of order two.
The invariant integral kernel is primitive, so such a
primitive vector is obtained by Smith reduction.

Finite order is also effectively bounded at this linear
stage. Its minimal polynomial is a product of distinct
cyclotomic polynomials, and eigenvalue one uses at least one
of the six dimensions. Every other cyclotomic factor therefore
has degree at most five and their total degree is at most
five. For a cyclotomic index $d$, each prime divisor $p$ has
$p-1\le5$, and its prime-power factors satisfy
$p^{a-1}(p-1)\le5$. Enumerating these factors and the possible
degree sums gives exactly
\[
                 d\in\{1,2,3,4,5,6,8,10,12\}.             \tag{F.CN1}
\]
Here $d$ is the order of the matrix, the least common multiple
of its cyclotomic indices. Multiplicities of a factor can only
increase the dimension and introduce no new orders. Each
listed order occurs with a fixed vector after adding trivial
summands. Thus taking the powers in (F.CN1), checking the
fixed kernel, the product relation, and the cusp Smith
factors gives an exact admissibility test. The unbounded
actual affine multiplicities are not bounded by (F.CN1).

For an admissible full tuple $T=(M_1,\ldots,M_k,M_0)$ let
$H(T)$ be its maximum absolute entry. If a three-point linear
tuple has an identity at a finite position, temporarily delete
that identity, classify the remaining two-point tuple, and
reinsert the identity at the same marked position. There can
be only one such identity, since otherwise a finite matrix
would equal the inverse of the nontrivial unipotent cusp.
Deleting a fixed identity changes neither height nor the
lexicographic comparison of tuples in that conjugacy-invariant
sector. This is only a linear normalization: a nontrivial
affine order at the identity-linear site is retained in the
gluing datum. In the two-involution sector use the canonical
representative defined in F.2.38. This sector is invariant under simultaneous
conjugacy. In every other sector, order full integer tuples
first by $H$, then lexicographically by their row-major
entries in the displayed marked order. Define the canonical
representative to be the first admissible tuple in the
simultaneous conjugacy class.

To compute it, enumerate heights $1,\ldots,H(T)$ and the
integer entries at that height. There are at most
\[
                     (2H(T)+1)^{36(k+1)}                  \tag{F.CN2}
\]
tuples in this complete box. The product relation allows
$M_0$ to be calculated from the first $k$ matrices, so the
implementation enumerates fewer entries. It also rejects
matrices with different characteristic polynomials or orders
before calling the conjugacy decision F.2.36. These filters
cannot reject an equivalent tuple. At least one comparison
succeeds: the supplied tuple itself belongs to its own box.
The first successful comparison returns the canonical
representative and its exact conjugating matrix.

Every simultaneous orbit has a least element in this order.
Its height is at most the height of any supplied member,
which proves that (F.CN2) loses no class representative for
that input. If two inputs are equivalent they therefore
return the same least element. Conversely the returned
conjugators show that equal representatives imply equivalent
inputs. The casewise involution normal form has the same
property by F.2.38, and an involution tuple cannot be
equivalent to a tuple in the other branch. This proves the
classification assertion, including its finite stopping rule.

Finally, F.2.16 already gives an unbounded admissible integral
family. In its homology convention direct symbolic
multiplication gives
\[
             \operatorname{tr}(M_1M_2^3)
                       =-v^4+v^3-v^2+2.                 \tag{F.CN3}
\]
This is unbounded even on the previously constructed free
affine subfamily $v\in20\mathbb Z$. It is invariant under
simultaneous conjugacy. A fixed finite entry box contains
only finitely many tuples and hence only finitely many such
traces; it cannot contain a representative of every class.
The input-dependent nature of (F.CN2) is therefore necessary
for this particular exhaustive procedure.
\end{proof}

The public implementation is
\path{code/fourfolds_tuple_classification.py}.
Its finite involution branch has been run on all $56$
representatives after nontrivial simultaneous integral
conjugation, and returned every exact representative and
conjugator. F.2.38 separately records $112$ normalization
checks. The marked identity-insertion branch was also run
on all $168$ conjugated inputs from the 56 classes and the
three possible identity positions; every representative and
conjugator passed. The general height-box fallback has not been
exhaustively run on a noninvolution tuple; no such computation
is claimed. Its underlying exact comparison procedure was
run on the twelve instances listed in F.2.36, including
three actual order-$(5,4)$ classes. The symbolic trace
calculation is saved with its executable program and output.
The all-twist group algorithm is F.2.31, and the full
necessary-and-sufficient three-connected-survivor tests are
F.2.30. These uses of the classifier do not presume that
arbitrarily many classes can be tabulated in finite time.

When transporting a specified glued manifold, the returned
matrix $P$ also sends each affine numerator to $Pv_i$, the
clutch vector to $P\kappa$, and the supplied integral cusp
marking to $P P_{\mathrm{cusp}}$. The multiplicities are
unchanged. These are the actual linear changes of fibre
coordinates and therefore carry every boundary gluing map.
The linear canonical representative does not authorize
discarding the cusp marking and independently choosing a
different Smith marking for the filled manifold.
